\documentclass[10pt,a4paper]{amsart}
\usepackage{amsmath,amssymb,amsthm,graphicx}
% Retained preamble for the Stacks Project chapter {CH} (shared project preamble).
% The chapter file pinned-original-source.tex only carries \input{{preamble}} on its first line,
% so no preamble text lives inside the chapter itself.  The Stacks Project uses one single shared
% preamble for every chapter: the snapshot stages it three times (sources/stacks-chow/preamble.tex,
% sources/stacks-etale-cohomology/preamble.tex and focused-seeds/AGP/sources/stacks-snapshot-20261011/
% preamble.tex) and those copies are identical in size.  The block below is copied mechanically,
% line for line, from sources/stacks-chow/preamble.tex lines 14-18 and 154-217: the Xy-pic setup and
% the native theorem declarations together with the mathematical macros.  Nothing was added, removed
% or redefined; the wrapper class and the amsmath/amssymb/amsthm/graphicx base packages are supplied
% by the assembler.
% For commutative diagrams we use Xy-pic
% Chapter: stacks-cohomology.tex -- Cohomology of Sheaves.
\usepackage[all]{xy}

% We use 2cell for 2-commutative diagrams.
\xyoption{2cell}

% Theorem environments.
%
\theoremstyle{plain}
\newtheorem{theorem}[subsection]{Theorem}
\newtheorem{proposition}[subsection]{Proposition}
\newtheorem{lemma}[subsection]{Lemma}

\theoremstyle{definition}
\newtheorem{definition}[subsection]{Definition}
\newtheorem{example}[subsection]{Example}
\newtheorem{exercise}[subsection]{Exercise}
\newtheorem{situation}[subsection]{Situation}

\theoremstyle{remark}
\newtheorem{remark}[subsection]{Remark}
\newtheorem{remarks}[subsection]{Remarks}

\numberwithin{equation}{subsection}

% Macros
%
\def\lim{\mathop{\mathrm{lim}}\nolimits}
\def\colim{\mathop{\mathrm{colim}}\nolimits}
\def\Spec{\mathop{\mathrm{Spec}}}
\def\Hom{\mathop{\mathrm{Hom}}\nolimits}
\def\Ext{\mathop{\mathrm{Ext}}\nolimits}
\def\SheafHom{\mathop{\mathcal{H}\!\mathit{om}}\nolimits}
\def\SheafExt{\mathop{\mathcal{E}\!\mathit{xt}}\nolimits}
\def\Sch{\mathit{Sch}}
\def\Mor{\mathop{\mathrm{Mor}}\nolimits}
\def\Ob{\mathop{\mathrm{Ob}}\nolimits}
\def\Sh{\mathop{\mathit{Sh}}\nolimits}
\def\NL{\mathop{N\!L}\nolimits}
\def\CH{\mathop{\mathrm{CH}}\nolimits}
\def\proetale{{pro\text{-}\acute{e}tale}}
\def\etale{{\acute{e}tale}}
\def\QCoh{\mathit{QCoh}}
\def\Ker{\mathop{\mathrm{Ker}}}
\def\Im{\mathop{\mathrm{Im}}}
\def\Coker{\mathop{\mathrm{Coker}}}
\def\Coim{\mathop{\mathrm{Coim}}}

% Boxtimes
%
\DeclareMathSymbol{\boxtimes}{\mathbin}{AMSa}{"02}

%
% Macros for moduli stacks/spaces
%
\def\QCohstack{\mathcal{QC}\!\mathit{oh}}
\def\Cohstack{\mathcal{C}\!\mathit{oh}}
\def\Spacesstack{\mathcal{S}\!\mathit{paces}}
\def\Quotfunctor{\mathrm{Quot}}
\def\Hilbfunctor{\mathrm{Hilb}}
\def\Curvesstack{\mathcal{C}\!\mathit{urves}}
\def\Polarizedstack{\mathcal{P}\!\mathit{olarized}}
\def\Complexesstack{\mathcal{C}\!\mathit{omplexes}}
% \Pic is the operator that assigns to X its picard group, usage \Pic(X)
% \Picardstack_{X/B} denotes the Picard stack of X over B
% \Picardfunctor_{X/B} denotes the Picard functor of X over B
\def\Pic{\mathop{\mathrm{Pic}}\nolimits}
\def\Picardstack{\mathcal{P}\!\mathit{ic}}
\def\Picardfunctor{\mathrm{Pic}}
\def\Deformationcategory{\mathcal{D}\!\mathit{ef}}

% --- appended 2026-10-12 (Chief): the project own definitions of the reference/slogan/history
% display environments, copied mechanically from the shared Stacks preamble lines 8-12
% (\usepackage{verbatim} plus the three \newenvironment lines, which the chapter preamble keeps
% verbatim from the project preamble). Nothing added, removed or redefined: these environments
% carry no mathematical content and the project itself typesets them as comments. Without these
% three lines a counted statement whose own text uses them fails with Environment slogan undefined.
\usepackage{verbatim}
\newenvironment{reference}{\comment}{\endcomment}
\newenvironment{slogan}{\comment}{\endcomment}
\newenvironment{history}{\comment}{\endcomment}
\usepackage{cleveref}
\crefname{originalreference}{original reference}{original references}
\Crefname{originalreference}{Original reference}{Original references}
\makeatletter
\newlabel{more-algebra-example-spectral-sequence-principal}{{0BKG}{source}{}{sourcepointer-1}{}}
\newlabel{more-algebra-example-spectral-sequence-principal@cref}{{[originalreference][0][]0BKG}{[1][0][]source}{}{}{}}
\makeatother
\crefname{originallemma}{original lemma}{original lemmas}
\Crefname{originallemma}{Original lemma}{Original lemmas}
\makeatletter
\newlabel{lemma-Rf-commutes-with-Rlim}{{0BKP}{source}{}{sourcepointer-2}{}}
\newlabel{lemma-Rf-commutes-with-Rlim@cref}{{[originallemma][0][]0BKP}{[1][0][]source}{}{}{}}
\makeatother
\crefname{originalreference}{original reference}{original references}
\Crefname{originalreference}{Original reference}{Original references}
\makeatletter
\newlabel{more-algebra-lemma-compare-completion-derived-completion}{{0H32}{source}{}{sourcepointer-3}{}}
\newlabel{more-algebra-lemma-compare-completion-derived-completion@cref}{{[originalreference][0][]0H32}{[1][0][]source}{}{}{}}
\makeatother
\title{Formal functions for a principal ideal: derived completion of RΓ}
\author{The Stacks Project Authors}
\date{stacks-project:cohomology.tex:local-snapshot-20261011; GNU FDL 1.2 or later}
\begin{document}
\maketitle
\noindent\textit{Editorial note.} All delivered mathematical text is the author's own, copied verbatim from the pinned Stacks Project chapter file cohomology.tex; no statement or proof step was rewritten, reordered or removed.\par\medskip
\noindent\textit{Editorial note.} The counted span cites the lemmas recorded in assembly.external\_references with canonical tags as navigation pointers.\par\medskip
\noindent\textit{Allowed standard prerequisites (editorial).}\begin{itemize}
\item RΓ commutes with R lim. Cohomology of Sheaves, Lemma cohomology-lemma-Rf-commutes-with-Rlim. RΓ(X, R lim E\_n) = R lim RΓ(X,E\_n).
\item Completion vs derived completion. More on Algebra, Lemma more-algebra-lemma-compare-completion-derived-completion. The commutative diagram with the f-torsion T\_f.
\end{itemize}\medskip
\ifdefined\hypertarget\hypertarget{sourcepointer-1}{}\fi
\noindent\textit{Source reference.} Original reference 0BKG: sources/stacks-chow/tags.txt, source lines 14559--14559. Cross-chapter standard prerequisite cited by the author inside the counted span; it is used as a pointer and its own proof is not part of this item.\par\medskip
\ifdefined\hypertarget\hypertarget{sourcepointer-2}{}\fi
\noindent\textit{Source reference.} Original lemma 0BKP: sources/stacks-chow/tags.txt, source lines 14567--14567. Cross-chapter standard prerequisite cited by the author inside the counted span; it is used as a pointer and its own proof is not part of this item.\par\medskip
\ifdefined\hypertarget\hypertarget{sourcepointer-3}{}\fi
\noindent\textit{Source reference.} Original reference 0H32: sources/stacks-chow/tags.txt, source lines 21302--21302. Cross-chapter standard prerequisite cited by the author inside the counted span; it is used as a pointer and its own proof is not part of this item.\par\medskip
\begin{lemma}
\label{lemma-formal-functions-principal}
Let $(X, \mathcal{O}_X)$ be a ringed space. Let $A \to \Gamma(X, \mathcal{O}_X)$
be a ring map and let $f \in A$. Let $E$ be an object of $D(\mathcal{O}_X)$.
Denote
$$
E_n = E \otimes_{\mathcal{O}_X} (\mathcal{O}_X \xrightarrow{f^n} \mathcal{O}_X)
$$
and set $E^\wedge = R\lim E_n$. For $p \in \mathbf{Z}$
is a canonical commutative diagram
$$
\xymatrix{
& 0 & 0 \\
0 \ar[r] &
\widehat{H^p(X, E)} \ar[r] \ar[u] &
\lim H^p(X, E_n) \ar[r] \ar[u] &
T_f(H^{p + 1}(X, E)) \ar[r] &
0 \\
0 \ar[r] &
H^0(H^p(X, E)^\wedge) \ar[r] \ar[u] &
H^p(X, E^\wedge) \ar[r] \ar[u] &
T_f(H^{p + 1}(X, E)) \ar[r] \ar@{=}[u] &
0 \\
&
R^1\lim H^p(X, E)[f^n] \ar[u] \ar[r]^\cong &
R^1\lim H^{p - 1}(X, E_n) \ar[u] \\
& 0 \ar[u] & 0 \ar[u]
}
$$
with exact rows and columns
where $\widehat{H^p(X, E)} = \lim H^p(X, E)/f^n H^p(X, E)$ is the usual
$f$-adic completion, $H^p(X, E)^\wedge$ is the derived $f$-adic completion,
and $T_f(H^{p + 1}(X, E))$ is the $f$-adic Tate module, see
More on Algebra, Example
\ref{more-algebra-example-spectral-sequence-principal}.
Finally, we have $H^p(X, E^\wedge) = H^p(R\Gamma(X, E)^\wedge)$.
\end{lemma}

\begin{proof}
Observe that $R\Gamma(X, E^\wedge) = R\lim R\Gamma(X, E_n)$ by
Lemma \ref{lemma-Rf-commutes-with-Rlim}.
On the other hand, we have
$$
R\Gamma(X, E_n) =
R\Gamma(X, E) \otimes_A^\mathbf{L} (A \xrightarrow{f^n} A)
$$
(details omitted). We find that $R\Gamma(X, E^\wedge)$ is the derived
$f$-adic completion $R\Gamma(X, E)^\wedge$. Whence the diagram by
More on Algebra, Lemma
\ref{more-algebra-lemma-compare-completion-derived-completion}.
\end{proof}

\end{document}
